\subsubsection{The Hahn-Banach Theorem}\label{subsec:hahn-banach-theorem}

In the last example of the preceding section, we used the following result.

\begin{Theorem}
\label{thm:hahn-banach-extension}
Let $X$ be a Banach space and let $M$ be a linear subspace of $X$. Let $l_M$ be a bounded linear mapping from $M$ to $\mathbb{R}$ ($\mathbb{C}$). Then there is a linear functional $l\in X^*$ such that $l|_M=l_M$ and $\|l\|=\|l_M\|$.
\end{Theorem}

In other words, linear functionals defined on any linear subspace can be extended without increasing their norm. This result and various more general ones from which it follows are known as the Hahn-Banach theorem. Some of the more general versions are relevant in convex analysis and optimization; since we shall not pursue these topics in the book, we do not state those results. The Hahn-Banach theorem guarantees the existence of many linear functionals, because we can always define linear functionals on any finite-dimensional subspace and then extend them. In particular, given any nonzero $x\in X$, there exists $f\in X^*$ with $f(x)\ne0$.

For the proof of Theorem~\ref{thm:hahn-banach-extension}, we refer the reader to texts on functional analysis. For the case of a Hilbert space, however, the result is almost trivial. We first extend $l$ to the closure of $M$ by continuity, set $l=0$ on $M^\perp$ and then define $l$ on the whole space by linearity. We leave it to the reader to verify that the $l$ defined in this way has the properties claimed in the theorem; see Problem~\ref{ex:hahn-banach-hilbert-extension-details}.
